\section{Introduction}

Four-phase polarization in the V59 compiler places the phases in the literal source
sequences $\beta+\ii^jw$ and uses one common Fourier--Poisson profile
\cite{tpc59}.  TPC-227 proved that a packet-dependent transform preserves this compiler
for all sources only when its Gram equals the physical common Gram \cite{tpc227}.
That theorem removes a typing ambiguity but leaves a constructive question: what exact
collision scalar results when the source phase is kept in the right place?

We answer that question at the finite Hilbert level.  Prime-labelled row transforms of
$\beta$ and $w$ are denoted $U_q$ and $V_q$.  The natural packet row is
$U_q+\ii^jV_q$.  Instead of taking an unsigned envelope, we first subtract the
same-prime quadratic diagonal and then polarize.  The result is exactly the ordered
off-diagonal source correlation
\[
 \sum_{q\ne r}\langle U_q,V_r\rangle.
\]

The first primitive collision classified in TPC-226 \cite{tpc226} then supplies a
minimal four-term block.  This does not determine an arithmetic sign, but it replaces
the former free-profile sign question by a literal source-labelled target.
